\documentclass[12pt,a4paper]{article}
\usepackage[utf8]{inputenc}
\usepackage[T1]{fontenc}
\usepackage[brazil]{babel}
\usepackage{helvet}
\renewcommand{\familydefault}{\sfdefault}
\usepackage{geometry,setspace,longtable,array,indentfirst,hyperref,enumitem,ragged2e}
\usepackage[table]{xcolor}
\usepackage{titlesec,fancyhdr}

\definecolor{AzulCorporativo}{RGB}{0,82,155}
\definecolor{AzulPetroleo}{RGB}{23,104,133}
\definecolor{CinzaEscuro}{RGB}{64,64,64}
\definecolor{LaranjaDestaque}{RGB}{230,115,0}
\hypersetup{colorlinks=true,linkcolor=AzulPetroleo,urlcolor=AzulCorporativo,
	citecolor=LaranjaDestaque,
	pdftitle={Modelagem Lógica - Modelo Analítico de Dados para E-commerce},
	pdfauthor={Lucas Dias Noronha}}
\titleformat{\section}{\color{AzulCorporativo}\normalfont\Large\bfseries}{\thesection}{1em}{}
\titleformat{\subsection}{\color{AzulPetroleo}\normalfont\large\bfseries}{\thesubsection}{1em}{}
\titleformat{\subsubsection}{\color{CinzaEscuro}\normalfont\normalsize\bfseries}{\thesubsubsection}{1em}{}
\geometry{left=3cm,right=2cm,top=3cm,bottom=2cm,headheight=45pt}
\onehalfspacing
\setlength{\parskip}{0.35em}
\pagestyle{fancy}
\fancyhf{}
\renewcommand{\headrulewidth}{0pt}
\fancyhead[C]{\colorbox{AzulCorporativo}{\makebox[\dimexpr\textwidth-2\fboxsep\relax][l]{\textcolor{white}{\textbf{\ \ Modelagem Lógica}}}}}
\fancyfoot[C]{\color{CinzaEscuro}\thepage}
\newcommand{\projeto}{Modelo Analítico de Dados para E-commerce}
\newcommand{\autor}{Lucas Dias Noronha}
\newcommand{\dataDocumento}{\today}

\begin{document}
	\thispagestyle{empty}
	\vspace*{\fill}
	\begin{center}
		{\Huge \textcolor{AzulCorporativo}{\textbf{Modelagem Lógica}}}\\[0.5cm]
		{\Large \textcolor{AzulPetroleo}{\projeto}}\\[0.8cm]
	\end{center}
	\vspace*{\fill}
	\noindent
	\textbf{\textcolor{AzulCorporativo}{Autor:}} \autor
	\hfill
	\textbf{\textcolor{AzulCorporativo}{Data:}} \dataDocumento
	\newpage
	\setcounter{tocdepth}{1}
	\tableofcontents
	\newpage

	\section{Introdução}
	Este documento registra o desenvolvimento incremental do modelo lógico
	relacional do projeto \textit{\projeto}, construído a partir do modelo
	conceitual.

	O modelo lógico converte as nove entidades conceituais em tabelas, define suas
	chaves e relacionamentos, consolida a nomenclatura SQL, avalia a normalização
	até a Terceira Forma Normal e valida a integridade referencial. O resultado
	preserva semântica, granularidade e rastreabilidade e permanece independente de
	decisões físicas do Sistema Gerenciador de Banco de Dados.

	\section{Escopo da Etapa}
	Esta etapa contempla:
	\begin{itemize}[leftmargin=1.2cm]
		\item a correspondência entre entidades conceituais e tabelas lógicas;
		\item a conversão dos atributos próprios das entidades em colunas;
		\item a preservação da granularidade aprovada no modelo conceitual;
		\item a definição de chaves primárias, chaves estrangeiras, nulabilidade e
		restrições lógicas de unicidade;
		\item a aplicação da convenção de nomenclatura SQL;
		\item a avaliação da normalização até a Terceira Forma Normal;
		\item a validação da integridade referencial e a revisão de redundâncias;
		\item o registro de decisões, divergências identificadas e implicações para a futura modelagem física.
	\end{itemize}

	Permanecem fora do escopo os tipos específicos do PostgreSQL, índices físicos,
	particionamento, estratégias de armazenamento e scripts DDL. Essas decisões
	pertencem à modelagem física.

	\section{Conversão de Entidades em Tabelas}
	{\small
		\renewcommand{\arraystretch}{1.35}
		\begin{longtable}{|>{\raggedright\arraybackslash}p{3.2cm}|>{\raggedright\arraybackslash}p{3cm}|>{\raggedright\arraybackslash}p{8cm}|}
			\hline
			\rowcolor{AzulCorporativo}
			\textcolor{white}{\textbf{Entidade conceitual}} & \textcolor{white}{\textbf{Tabela lógica}} & \textcolor{white}{\textbf{Granularidade preservada}} \\ \hline
			\endfirsthead
			\hline
			\rowcolor{AzulCorporativo}
			\textcolor{white}{\textbf{Entidade conceitual}} & \textcolor{white}{\textbf{Tabela lógica}} & \textcolor{white}{\textbf{Granularidade preservada}} \\ \hline
			\endhead
			Cliente & \texttt{cliente} & Um registro de cliente associado a um pedido na origem. \\ \hline
			Pedido & \texttt{pedido} & Um pedido. \\ \hline
			Item do Pedido & \texttt{item\_pedido} & Um item sequencial dentro de um pedido. \\ \hline
			Produto & \texttt{produto} & Um produto. \\ \hline
			Vendedor & \texttt{vendedor} & Um vendedor. \\ \hline
			Pagamento & \texttt{pagamento} & Uma ocorrência sequencial de pagamento dentro de um pedido. \\ \hline
			Avaliação & \texttt{avaliacao} & Uma ocorrência de avaliação associada a um pedido. \\ \hline
			Prefixo CEP & \texttt{prefixo\_cep} & Uma referência territorial agregada por prefixo de CEP. \\ \hline
			Geolocalização & \texttt{geolocalizacao} & Uma observação geográfica associada a um prefixo de CEP. \\ \hline
	\end{longtable}}

	As nove entidades conceituais originaram nove tabelas lógicas. Não foi
	necessária a criação, fusão ou eliminação de tabelas nesta etapa.

	\section{Conversão dos Atributos em Colunas}
	{\small
		\renewcommand{\arraystretch}{1.35}
		\begin{longtable}{|>{\raggedright\arraybackslash}p{3.1cm}|>{\raggedright\arraybackslash}p{11.3cm}|}
			\hline
			\rowcolor{AzulCorporativo}
			\textcolor{white}{\textbf{Tabela}} & \textcolor{white}{\textbf{Colunas provenientes do modelo conceitual}} \\ \hline
			\endfirsthead
			\hline
			\rowcolor{AzulCorporativo}
			\textcolor{white}{\textbf{Tabela}} & \textcolor{white}{\textbf{Colunas provenientes do modelo conceitual}} \\ \hline
			\endhead
			\texttt{cliente} & \texttt{id\_cliente}, \texttt{id\_cliente\_unico}, \texttt{cidade}, \texttt{estado}. \\ \hline
			\texttt{pedido} & \texttt{id\_pedido}, \texttt{status\_pedido}, \texttt{data\_compra}, \texttt{data\_aprovacao}, \texttt{data\_envio\_transportador}, \texttt{data\_entrega}, \texttt{data\_estimada}. \\ \hline
			\texttt{item\_pedido} & \texttt{id\_pedido}, \texttt{id\_item}, \texttt{preco\_item}, \texttt{valor\_frete}, \texttt{data\_limite\_envio}. \\ \hline
			\texttt{produto} & \texttt{id\_produto}, \texttt{nome\_categoria}, \texttt{comprimento\_nome}, \texttt{comprimento\_descricao}, \texttt{quantidade\_fotos}, \texttt{peso\_g}, \texttt{comprimento\_cm}, \texttt{altura\_cm}, \texttt{largura\_cm}. \\ \hline
			\texttt{vendedor} & \texttt{id\_vendedor}, \texttt{cidade}, \texttt{estado}. \\ \hline
			\texttt{pagamento} & \texttt{id\_pedido}, \texttt{sequencial\_pagamento}, \texttt{tipo\_pagamento}, \texttt{numero\_parcelas}, \texttt{valor\_pagamento}. \\ \hline
			\texttt{avaliacao} & \texttt{id\_avaliacao}, \texttt{id\_pedido}, \texttt{nota\_avaliacao}, \texttt{data\_criacao}, \texttt{data\_resposta}. \\ \hline
			\texttt{prefixo\_cep} & \texttt{prefixo\_cep}. \\ \hline
			\texttt{geolocalizacao} & \texttt{id\_geolocalizacao}, \texttt{latitude}, \texttt{longitude}, \texttt{cidade}, \texttt{estado}. \\ \hline
	\end{longtable}}

	\section{Decisões e Limites da Conversão}
	\subsection{Atributos Próprios e Relacionamentos}
	O MER conceitual representa relacionamentos como associações, e não como
	atributos duplicados nas entidades participantes. Por isso, durante a conversão
	inicial das entidades em tabelas não foram introduzidos antecipadamente:
	\begin{itemize}[leftmargin=1.2cm]
		\item \texttt{id\_cliente} em \texttt{pedido};
		\item \texttt{id\_produto} e \texttt{id\_vendedor} em \texttt{item\_pedido};
		\item \texttt{prefixo\_cep} em \texttt{cliente}, \texttt{vendedor} e \texttt{geolocalizacao}.
	\end{itemize}

	Durante a conversão, essas colunas foram materializadas como chaves estrangeiras. O
	atributo \texttt{id\_pedido}, já presente nas tabelas \texttt{item\_pedido},
	\texttt{pagamento} e \texttt{avaliacao}, foi mantido por integrar seus
	identificadores conceituais.

	\subsection{Prefixo CEP}
	A análise de requisitos original descrevia oito entidades e armazenava o
	prefixo de CEP diretamente nas estruturas geográficas. Durante a modelagem conceitual, Prefixo CEP foi promovido a uma nona entidade, separando a
	referência territorial agregada das múltiplas observações de geolocalização.
	Esta conversão preserva o refinamento aprovado.

	\subsection{Avaliação}
	A relação Pedido--Avaliação foi inicialmente considerada 1:1 nos requisitos,
	mas a validação empírica realizada durante a modelagem conceitual estabeleceu cardinalidade 1:N. A tabela
	\texttt{avaliacao} mantém a granularidade de uma ocorrência identificada pelo
	par conceitual \texttt{id\_avaliacao} e \texttt{id\_pedido}, posteriormente
	preservado como chave primária lógica composta.

	\subsection{Cliente e Pedido}
	O modelo conceitual consolidado registra Cliente--Pedido como 1:1 porque o
	identificador de cliente distingue a ocorrência associada ao pedido na origem. O
	atributo longitudinal do consumidor pode repetir-se entre ocorrências. Nenhuma consolidação foi realizada.

	\subsection{Domínios Lógicos}
	O esquema técnico utiliza domínios descritivos, como \texttt{identificador},
	\texttt{data\_hora}, \texttt{valor\_monetario} e
	\texttt{coordenada\_geografica}. Eles comunicam a natureza lógica das colunas
	sem escolher tipos ou recursos específicos do PostgreSQL. A definição física
	permanece reservada à futura modelagem física.

	No artefato DBML, o bloco \texttt{indexes} declara propriedades lógicas de
	chave primária e unicidade exigidas pela sintaxe da ferramenta. Ele não
	representa, nesta etapa, uma estratégia de indexação física.

	\section{Definição das Chaves Primárias Lógicas}

	\subsection{Critérios Adotados}

	As chaves foram avaliadas quanto a unicidade, completude, estabilidade,
	granularidade e aderência ao identificador conceitual. Chaves compostas foram
	mantidas quando necessárias para distinguir ocorrências dentro do contexto do
	pedido. A chave substituta de Geolocalização, definida durante a modelagem conceitual, deverá ser
	gerada de forma estável durante a carga e não possui equivalente direto na
	origem.

	{\small
		\renewcommand{\arraystretch}{1.4}
		\begin{longtable}{|>{\raggedright\arraybackslash}p{2.7cm}|>{\raggedright\arraybackslash}p{5.8cm}|>{\raggedright\arraybackslash}p{2.2cm}|>{\raggedright\arraybackslash}p{3cm}|}
			\hline
			\rowcolor{AzulCorporativo}
			\textcolor{white}{\textbf{Tabela}} & \textcolor{white}{\textbf{Chave primária lógica}} & \textcolor{white}{\textbf{Estrutura}} & \textcolor{white}{\textbf{Restrição}} \\ \hline
			\endfirsthead
			\hline
			\rowcolor{AzulCorporativo}
			\textcolor{white}{\textbf{Tabela}} & \textcolor{white}{\textbf{Chave primária lógica}} & \textcolor{white}{\textbf{Estrutura}} & \textcolor{white}{\textbf{Restrição}} \\ \hline
			\endhead
			\texttt{cliente} & \texttt{id\_cliente} & Natural simples & \texttt{pk\_cliente} \\ \hline
			\texttt{pedido} & \texttt{id\_pedido} & Natural simples & \texttt{pk\_pedido} \\ \hline
			\texttt{item\_pedido} & \texttt{id\_pedido}, \texttt{id\_item} & Natural composta & \texttt{pk\_item\_pedido} \\ \hline
			\texttt{produto} & \texttt{id\_produto} & Natural simples & \texttt{pk\_produto} \\ \hline
			\texttt{vendedor} & \texttt{id\_vendedor} & Natural simples & \texttt{pk\_vendedor} \\ \hline
			\texttt{pagamento} & \texttt{id\_pedido}, \texttt{sequencial\_pagamento} & Natural composta & \texttt{pk\_pagamento} \\ \hline
			\texttt{avaliacao} & \texttt{id\_avaliacao}, \texttt{id\_pedido} & Natural composta & \texttt{pk\_avaliacao} \\ \hline
			\texttt{prefixo\_cep} & \texttt{prefixo\_cep} & Natural simples & \texttt{pk\_prefixo\_cep} \\ \hline
			\texttt{geolocalizacao} & \texttt{id\_geolocalizacao} & Substituta simples & \texttt{pk\_geolocalizacao} \\ \hline
	\end{longtable}}

	A validação nos CSVs confirmou ausência de nulos e duplicidades nas chaves. Em
	Avaliação, os 99.224 pares formados pelos identificadores da avaliação e do
	pedido são únicos. O identificador de avaliação isolado possui 98.410 valores
	distintos e não é único na fonte.

	\section{Definição das Chaves Estrangeiras}

	Os relacionamentos do MER foram materializados por chaves estrangeiras. A
	participação de cada tabela dependente é obrigatória; por isso, as colunas de
	referência são logicamente não nulas. Exclusões em cascata não foram adotadas:
	a política lógica é restringir a exclusão de registros referenciados e não
	propagar atualizações de identificadores.

	{\small
		\renewcommand{\arraystretch}{1.4}
		\begin{longtable}{|>{\raggedright\arraybackslash}p{4.6cm}|>{\raggedright\arraybackslash}p{5.2cm}|>{\raggedright\arraybackslash}p{4.8cm}|}
			\hline
			\rowcolor{AzulCorporativo}
			\textcolor{white}{\textbf{Restrição}} & \textcolor{white}{\textbf{Origem}} & \textcolor{white}{\textbf{Destino}} \\ \hline
			\endfirsthead
			\hline
			\rowcolor{AzulCorporativo}
			\textcolor{white}{\textbf{Restrição}} & \textcolor{white}{\textbf{Origem}} & \textcolor{white}{\textbf{Destino}} \\ \hline
			\endhead
			\texttt{fk\_cliente\_prefixo\_cep} & \texttt{cliente.prefixo\_cep} & \texttt{prefixo\_cep.prefixo\_cep} \\ \hline
			\texttt{fk\_pedido\_cliente} & \texttt{pedido.id\_cliente} & \texttt{cliente.id\_cliente} \\ \hline
			\texttt{fk\_item\_pedido\_pedido} & \texttt{item\_pedido.id\_pedido} & \texttt{pedido.id\_pedido} \\ \hline
			\texttt{fk\_item\_pedido\_produto} & \texttt{item\_pedido.id\_produto} & \texttt{produto.id\_produto} \\ \hline
			\texttt{fk\_item\_pedido\_vendedor} & \texttt{item\_pedido.id\_vendedor} & \texttt{vendedor.id\_vendedor} \\ \hline
			\texttt{fk\_vendedor\_prefixo\_cep} & \texttt{vendedor.prefixo\_cep} & \texttt{prefixo\_cep.prefixo\_cep} \\ \hline
			\texttt{fk\_pagamento\_pedido} & \texttt{pagamento.id\_pedido} & \texttt{pedido.id\_pedido} \\ \hline
			\texttt{fk\_avaliacao\_pedido} & \texttt{avaliacao.id\_pedido} & \texttt{pedido.id\_pedido} \\ \hline
			\texttt{fk\_geolocalizacao\_\allowbreak prefixo\_cep} & \texttt{geolocalizacao.prefixo\_cep} & \texttt{prefixo\_cep.prefixo\_cep} \\ \hline
	\end{longtable}}

	A relação Cliente--Pedido é materializada pela chave estrangeira obrigatória
	\texttt{pedido.id\_\hspace{0pt}cliente} e pela restrição de unicidade
	\texttt{uq\_pedido\_id\_\hspace{0pt}cliente}. Em conjunto, essas regras garantem que todo
	Pedido referencie exatamente um Cliente e que cada Cliente seja referenciado
	por, no máximo, um Pedido. Elas não garantem, isoladamente, que todo Cliente
	possua um Pedido; a obrigatoriedade no sentido Cliente $\rightarrow$ Pedido
	depende do processo de carga ou de validação adicional. Assim, o modelo
	preserva o 1:1 observado nos registros carregados sem atribuir às restrições uma
	garantia bidirecional que elas não impõem por si sós.

	A cardinalidade preservada é Pedido 1:N Avaliação. Um pedido pode possuir zero
	ou mais avaliações, enquanto cada ocorrência de avaliação pertence
	obrigatoriamente a exatamente um pedido.

	\section{Convenção de Nomenclatura SQL}

	A convenção lógica utiliza caracteres ASCII e nomes em português, no singular.
	A separação entre palavras segue o padrão \texttt{snake\_case}. Tabelas e
	colunas não recebem prefixos técnicos nem abreviações não consolidadas.

	\begin{itemize}[leftmargin=1.2cm]
		\item chaves primárias: \texttt{pk\_}\textit{tabela};
		\item chaves estrangeiras: \texttt{fk\_}\textit{origem}\texttt{\_}\textit{destino};
		\item restrições de unicidade: \texttt{uq\_}\textit{tabela}\texttt{\_}\textit{atributos};
		\item verificações de domínio: \texttt{ck\_}\textit{tabela}\texttt{\_}\textit{regra};
		\item índices físicos: reservados à modelagem física, no padrão \texttt{ix\_}\textit{tabela}\texttt{\_}\textit{atributos}.
	\end{itemize}

	A convenção também qualifica nomes potencialmente ambíguos. Foram adotados
	\texttt{status\_pedido}, \texttt{preco\_item}, \texttt{nome\_categoria} e
	\texttt{nota\_avaliacao}. A acentuação permanece apenas na documentação. Nos
	identificadores relacionais são usados \texttt{avaliacao} e
	\texttt{geolocalizacao}.

	\section{Normalização até a Terceira Forma Normal}

	\subsection{Dependências Funcionais}

	A avaliação considerou atributos atômicos, dependência da chave inteira e
	ausência de dependências transitivas entre atributos não chave.

	{\small
		\renewcommand{\arraystretch}{1.4}
		\begin{longtable}{|>{\raggedright\arraybackslash}p{2.7cm}|>{\raggedright\arraybackslash}p{8.4cm}|>{\raggedright\arraybackslash}p{3.0cm}|}
			\hline
			\rowcolor{AzulCorporativo}
			\textcolor{white}{\textbf{Tabela}} & \textcolor{white}{\textbf{Determinante dos atributos não chave}} & \textcolor{white}{\textbf{Resultado}} \\ \hline
			\endfirsthead
			\hline
			\rowcolor{AzulCorporativo}
			\textcolor{white}{\textbf{Tabela}} & \textcolor{white}{\textbf{Determinante dos atributos não chave}} & \textcolor{white}{\textbf{Resultado}} \\ \hline
			\endhead
			\texttt{cliente} & \texttt{id\_cliente} & 3FN \\ \hline
			\texttt{pedido} & \texttt{id\_pedido} & 3FN \\ \hline
			\texttt{item\_pedido} & \texttt{(id\_pedido, id\_item)} & 3FN \\ \hline
			\texttt{produto} & \texttt{id\_produto} & 3FN \\ \hline
			\texttt{vendedor} & \texttt{id\_vendedor} & 3FN \\ \hline
			\texttt{pagamento} & \texttt{(id\_pedido, sequencial\_pagamento)} & 3FN \\ \hline
			\texttt{avaliacao} & \texttt{(id\_avaliacao, id\_pedido)} & 3FN \\ \hline
			\texttt{prefixo\_cep} & Não possui atributos não chave & 3FN \\ \hline
			\texttt{geolocalizacao} & \texttt{id\_geolocalizacao} & 3FN \\ \hline
	\end{longtable}}

	\subsection{Identidade e Dependências de Avaliação}

	O atributo \texttt{id\_avaliacao} isoladamente não é único na fonte. A
	ocorrência de avaliação é identificada semanticamente pelo par
	\texttt{(id\_avaliacao, id\_pedido)}, conforme validado durante a modelagem conceitual. O segundo
	atributo também é uma chave estrangeira obrigatória para
	\texttt{pedido.id\_pedido}.

	A validação empírica encontrou 789 identificadores de avaliação associados a
	mais de um pedido. Nota, data de criação e data de resposta repetem-se nesses
	grupos. Essa coincidência é uma evidência empírica, mas não comprova a
	dependência funcional semântica segundo a qual \texttt{id\_avaliacao}
	determinaria \texttt{nota\_avaliacao}, \texttt{data\_criacao} e
	\texttt{data\_resposta}. A repetição será tratada como possível anomalia ou limitação de identidade do conjunto de dados.

	Como regra do domínio, os atributos descritivos da avaliação são considerados
	funcionalmente dependentes da chave composta completa. A chave composta não foi
	adotada para corrigir uma violação de normalização, mas para representar a
	unicidade observada na fonte. Sob essa interpretação semântica, não há violação
	reconhecida da 2FN, e \texttt{avaliacao} permanece em 2FN e 3FN. A decisão
	preserva a granularidade e a cardinalidade validadas durante a modelagem
	conceitual, sem alterá-la.

	\subsection{Atributos Geográficos}

	Não foi assumida a dependência \texttt{prefixo\_cep} $\rightarrow$
	\texttt{(cidade, estado)}. Foram encontrados prefixos associados a múltiplos
	pares de cidade e estado: 39 em Cliente, 49 em Vendedor e 8.559 em
	Geolocalização. Cidade e estado permanecem, portanto, atributos das ocorrências
	originais, evitando uma dependência funcional inexistente e uma consolidação
	indevida.

	\section{Validação da Integridade Referencial}

	A entidade \texttt{prefixo\_cep} é formada pela união dos prefixos presentes em
	Cliente, Vendedor e Geolocalização, totalizando 19.177 valores distintos. Essa
	estratégia preserva os 157 prefixos de clientes e os sete de vendedores que não
	possuem observações na fonte de geolocalização.

	{\small
		\renewcommand{\arraystretch}{1.25}
		\begin{longtable}{|>{\raggedright\arraybackslash}p{5.5cm}|>{\raggedright\arraybackslash}p{2.2cm}|>{\raggedright\arraybackslash}p{2.2cm}|>{\raggedright\arraybackslash}p{3.6cm}|}
			\hline
			\rowcolor{AzulCorporativo}
			\textcolor{white}{\textbf{Referência validada}} & \textcolor{white}{\textbf{Nulos}} & \textcolor{white}{\textbf{Órfãos}} & \textcolor{white}{\textbf{Conclusão}} \\ \hline
			\endfirsthead
			\hline
			\rowcolor{AzulCorporativo}
			\textcolor{white}{\textbf{Referência validada}} & \textcolor{white}{\textbf{Nulos}} & \textcolor{white}{\textbf{Órfãos}} & \textcolor{white}{\textbf{Conclusão}} \\ \hline
			\endhead
			Pedido $\rightarrow$ Cliente & 0 & 0 & Válida, obrigatória e única \\ \hline
			Item $\rightarrow$ Pedido & 0 & 0 & Válida e obrigatória \\ \hline
			Item $\rightarrow$ Produto & 0 & 0 & Válida e obrigatória \\ \hline
			Item $\rightarrow$ Vendedor & 0 & 0 & Válida e obrigatória \\ \hline
			Pagamento $\rightarrow$ Pedido & 0 & 0 & Válida e obrigatória \\ \hline
			Avaliação para Pedido & 0 & 0 & Válida e obrigatória \\ \hline
			Cliente $\rightarrow$ Prefixo CEP & 0 & 0 & Válida pela união \\ \hline
			Vendedor $\rightarrow$ Prefixo CEP & 0 & 0 & Válida pela união \\ \hline
			Geolocalização $\rightarrow$ Prefixo CEP & 0 & 0 & Válida pela união \\ \hline
	\end{longtable}}

	As verificações demonstram cobertura referencial integral. As regras lógicas
	\texttt{NOT NULL} aplicadas às FKs representam a participação obrigatória do
	lado dependente. Como política lógica, o DBML especifica
	\texttt{ON DELETE RESTRICT} e \texttt{ON UPDATE NO ACTION}: registros
	referenciados não são excluídos e alterações de identificadores não são
	propagadas automaticamente. A tradução dessas regras para os mecanismos
	específicos do PostgreSQL será confirmada durante a modelagem física.

	\section{Revisão de Redundâncias}

	A revisão final não identificou redundâncias indevidas remanescentes. Foram
	mantidas apenas repetições semanticamente necessárias:

	\begin{itemize}[leftmargin=1.2cm]
		\item cidade e estado em Cliente, Vendedor e Geolocalização, pois o prefixo de
		CEP não os determina univocamente nos dados;
		\item \texttt{id\_cliente\_unico} em Cliente, pois representa a identidade
		longitudinal do consumidor e não substitui a ocorrência \texttt{id\_cliente};
		\item identificadores propagados como FKs, necessários para materializar os
		relacionamentos e não considerados redundância descritiva.
	\end{itemize}

	A repetição de nota e datas para um mesmo identificador de avaliação foi
	mantida como possível anomalia ou limitação de identidade da fonte, pois não
	constitui uma dependência funcional semântica confirmada. A tradução de
	categorias continua fora do modelo central, conforme
	a decisão conceitual que a classificou como estrutura auxiliar de tratamento.
	Não foram introduzidos índices, particionamento ou desnormalização preventiva.

	\section{Modelo Lógico Consolidado}

	{\small
		\renewcommand{\arraystretch}{1.25}
		\begin{longtable}{|>{\raggedright\arraybackslash}p{3.2cm}|>{\raggedright\arraybackslash}p{4cm}|>{\raggedright\arraybackslash}p{3cm}|>{\raggedright\arraybackslash}p{3.7cm}|}
			\hline
			\rowcolor{AzulCorporativo}
			\textcolor{white}{\textbf{Tabela}} & \textcolor{white}{\textbf{PK}} & \textcolor{white}{\textbf{FKs}} & \textcolor{white}{\textbf{Demais atributos}} \\ \hline
			\endfirsthead
			\hline
			\rowcolor{AzulCorporativo}
			\textcolor{white}{\textbf{Tabela}} & \textcolor{white}{\textbf{PK}} & \textcolor{white}{\textbf{FKs}} & \textcolor{white}{\textbf{Demais atributos}} \\ \hline
			\endhead
			\texttt{cliente} & \texttt{id\_cliente} & \texttt{prefixo\_cep} & \texttt{id\_cliente\_unico}, \texttt{cidade}, \texttt{estado} \\ \hline
			\texttt{pedido} & \texttt{id\_pedido} & \texttt{id\_cliente} & \texttt{status\_pedido}, datas de processamento e entrega \\ \hline
			\texttt{item\_pedido} & \texttt{id\_pedido}, \texttt{id\_item} & \texttt{id\_pedido}, \texttt{id\_produto}, \texttt{id\_vendedor} & \texttt{data\_limite\_envio}, \texttt{preco\_item}, \texttt{valor\_frete} \\ \hline
			\texttt{produto} & \texttt{id\_produto} & --- & \texttt{nome\_categoria}, métricas textuais, fotos, peso e dimensões \\ \hline
			\texttt{vendedor} & \texttt{id\_vendedor} & \texttt{prefixo\_cep} & \texttt{cidade}, \texttt{estado} \\ \hline
			\texttt{pagamento} & \texttt{id\_pedido}, \texttt{sequencial\_pagamento} & \texttt{id\_pedido} & tipo, parcelas e valor \\ \hline
			\texttt{avaliacao} & \texttt{id\_avaliacao}, \texttt{id\_pedido} & \texttt{id\_pedido} & \texttt{nota\_avaliacao}, criação e resposta \\ \hline
			\texttt{prefixo\_cep} & \texttt{prefixo\_cep} & --- & --- \\ \hline
			\texttt{geolocalizacao} & \texttt{id\_geolocalizacao} & \texttt{prefixo\_cep} & latitude, longitude, cidade e estado \\ \hline
	\end{longtable}}

	O artefato técnico consolidado usa domínios lógicos, e não tipos físicos de
	PostgreSQL.
	\section{Prontidão para a Modelagem Física}

	O modelo consolidado possui nove tabelas em 3FN, chaves primárias e estrangeiras
	explícitas, nomenclatura uniforme e cobertura referencial validada. Permanecem
	corretamente reservadas à modelagem física as escolhas de tipos físicos e comprimentos,
	implementação da chave substituta de Geolocalização, restrições de domínio,
	índices e scripts DDL.

	\section{Resultado da Modelagem Lógica}

	O trabalho foi consolidado em um modelo relacional completo, rastreável e
	independente de SGBD. A tabela \texttt{avaliacao} preserva a chave composta e
	o relacionamento 1:N com Pedido validados durante a modelagem conceitual.
	Nenhuma alteração foi realizada no modelo conceitual.

\end{document}
